\section{Two complementary sensitivity references for a banking decision}
\label{exec23:sensitivity-reuse-scope}
A lending or servicing analysis can return a precise estimate while remaining
vulnerable to omitted borrower information. A useful product proposal must
explain which perturbation a sensitivity tool studies, which sampling
uncertainty it preserves, and what changes in the decision if the perturbation
is credible. Two concrete R candidates illustrate why one generic robustness
field would be inadequate. The inspected CRAN mirror revisions are
sensemakr 0.1.6, dated 23 July 2024, and EValue 4.1.4, dated 29 August 2025.
Their immutable trees authenticate thirteen selected source and documentation
blobs. They are different engines: the former develops regression sensitivity
in partial-$R^2$ coordinates; the latter includes risk-ratio sensitivity
thresholds and several effect-scale conversions.
\citep{exec23_sensitivity_sensemakr,exec23_sensitivity_evalue}

The package declarations are GPL-3 and GPL-2 respectively. The audit retains
those declarations and separately identifies local operating-system copies of
the license texts; it does not represent the latter as newly acquired package
release files or a review of every dependency. sensemakr declares R 3.1.0 or
later and suggests additional packages including fixest. EValue imports
several plotting, meta-analysis and data-processing dependencies. A named
maintainer and dated CRAN snapshot establish concrete provenance, but do not
promise future support or a complete reproducible environment. The reference
route below is therefore distinct from permission to distribute an adapter.
\citep{exec23_sensitivity_sensemakr,exec23_sensitivity_evalue}

\subsection{Regression sensitivity preserves its covariance convention}
\label{exec23:sensitivity-reuse-regression}
Let $\widehat\beta$ be the short-regression treatment coefficient, $s$ its
standard error and $\nu$ the residual degrees of freedom. Write
$a=R^2_{D\sim U\mid X}$ for omitted information $U$'s partial explanatory
strength for treatment $D$, conditional on recorded covariates $X$, and
$b=R^2_{Y\sim U\mid D,X}$ for its partial strength for outcome $Y$.
These are sensitivity restrictions, not estimates of an observed confounder.
The selected numeric helpers implement bias magnitude and adjusted standard
error
\[
 B=s\sqrt{\frac{\nu ab}{1-a}},\qquad
 s_{\mathrm{adj}}=s\sqrt{\frac{1-b}{1-a}\frac{\nu}{\nu-1}}.
\]
The sign and direction of adjustment are separate arguments: reducing a
positive coefficient gives $\widehat\beta-B$, while increasing its magnitude
is a different request. A credible banking report must identify the actual
regression, outcome units, adjustment set and sensitivity direction before
interpreting either number.
\citep{exec23_sensitivity_sensemakr}

For an entirely stipulated exact example, take $\widehat\beta=3$, $s=1$,
$\nu=8$, $a=1/5$ and $b=1/2$. Substitution gives $B=1$, adjusted coefficient
2 and adjusted variance $s_{\mathrm{adj}}^2=5/7$. Suppose one coefficient
unit is valued at USD 10 and the intervention costs USD 25. The original
point-based net is USD 5, whereas the sensitivity-adjusted point-based net
is USD $-5$. The sign change explains why plausible omitted-information
restrictions matter; it is not a measured treatment benefit, a confidence
statement or an empirical finding from either package.
\citep{exec23_sensitivity_reference}

In the inspected fixest path, model extraction explicitly requests iid
covariance. A warning in the corresponding adjustment path identifies this
restriction. Consequently, accepting a clustered model object does not justify
claiming that the original clustered uncertainty survives the sensitivity
calculation. The numeric helper's degrees-of-freedom guard admits zero and
one, while its adjusted-variance expression contains $\nu-1$ in the denominator.
The partial-$R^2$ validator admits the endpoint $a=1$, while the bias expression
contains $1-a$. These are static domain observations. The proposed regular
profile requires finite inputs, $s>0$, $\nu>1$, $0\leq a<1$ and $0\leq b<1$;
degenerate boundaries receive an explicit separate disposition instead of
being presented as a finite scientific result. No native failure is inferred
from source inspection alone.
\citep{exec23_sensitivity_sensemakr}

Two further details belong in the qualification protocol. The direct
\texttt{adjusted\_ci.numeric} helper uses $\nu-1$ in its critical-value
calculation; the inspected numeric wrapper constructs intervals with $\nu$.
Their uncertainty outputs cannot be promised identical without resolving the
intended convention. The wrapper also passes an estimate to an adjusted-SE
helper that accepts an ellipsis. That extra argument is accepted by the
signature; it is not evidence of an unused-argument error. Observed-covariate
benchmarking introduces additional assumptions through separate treatment and
outcome multipliers. Its implementation rejects impossible treatment bounds
and warns before capping an implied outcome bound above one. The adapter must
retain those warnings, raw restrictions and benchmark choices, rather than
turning a benchmark label into evidence that the omitted variable exists or
is no stronger than the chosen observed variable.
\citep{exec23_sensitivity_sensemakr}

\subsection{Risk-ratio sensitivity requires an explicit effect scale}
\label{exec23:sensitivity-reuse-ratio}
For the null target and a strictly positive observed risk ratio $r$, orient
its magnitude as $q=\max(r,1/r)$. The selected threshold calculation reduces to
\[
 E=q+\sqrt{q(q-1)}.
\]
The package describes this as confounding strength required to shift the
observed association to the specified target. It is not an estimated
confounder strength, a corrected treatment effect or a probability that the
causal claim is true. The interface also permits a non-null target, which must
remain an explicit field rather than silently inheriting the null formula.
It reports the relevant near-target confidence limit separately; a supplied
interval crossing the target receives threshold one. That result does not
establish that the effect is zero, and the original interval's sampling
justification remains the caller's responsibility.
\citep{exec23_sensitivity_evalue}

The reciprocal ratios $9/5$ and $5/9$ both give threshold 3, while the null
ratio 1 gives threshold 1. These exact arithmetic symmetries are useful
conversion oracles. They do not test finite-sample interval coverage. The
source rejects negative point ratios and inconsistent interval ordering, but
its zero-inclusive checks and reciprocal operations leave singular boundary
behavior for native qualification. The proposed worker instead declares a
strictly positive finite ratio and target, compatible bounds, event coding,
population and horizon; an infinite or missing result is not renamed strong
causal evidence.
\citep{exec23_sensitivity_evalue,exec23_sensitivity_reference}

The broader package offers different conversions for odds ratios, hazard
ratios, standardized mean differences and linear-regression coefficients.
For example, its nonrare-outcome odds-ratio conversion uses a square root.
This is an approximation, not a universal algebraic identity. For odds ratio
$\mathrm{OR}=81/25$, that conversion is $9/5$. Under a stipulated baseline
risk $p_0=1/5$, the exact risk-ratio identity for those two risks is instead
\[
 \mathrm{RR}=\frac{\mathrm{OR}}{1-p_0+p_0\mathrm{OR}}
             =\frac{405}{181},
\]
which differs from $9/5$. A banking interface that discards event prevalence
or the rare-outcome flag would conceal this distinction. The OLS route needs
an exposure contrast and a standard deviation; its documentation also states
that interval conversion treats the supplied standard deviation as known.
The conversion history and its additional assumptions must accompany the
returned sensitivity number.
\citep{exec23_sensitivity_evalue,exec23_sensitivity_reference}

\subsection{Reference decisions and the bounded prototype route}
\label{exec23:sensitivity-reuse-decision}
Retain both packages as method-specific proposal references. A future isolated
R worker is a conditional route for a declared regression-sensitivity or
risk-ratio profile, after complete artifact and dependency custody, applicable
rights review and native conformance. Direct adoption of every package family
would overextend this inspection. A new Python, Java or C++ port is not
justified without a measured platform or maintenance advantage sufficient to
cover dispatch, numerical, uncertainty and upgrade obligations. Independent
arithmetic remains a matched oracle, not a substitute calibrated inference
service. Neither package replaces graph identification, trend sensitivity,
selection-bias analysis or a different confounding model merely because all
are called sensitivity analysis.
\citep{exec23_sensitivity_sensemakr,exec23_sensitivity_evalue,exec23_sensitivity_reference}

The scientific owner must specify the causal question, adjustment clock,
sampling law and plausibility of sensitivity restrictions. The maintainer
must preserve model and row identity, covariance type, full conversion
history, warnings and failures; compare direct helpers with numeric and model
dispatch; and register singular-input and incompatible-interval cases. The
operator must bind an actual delivery, dependencies and resource budget.
Qualification then compares the selected outputs with independent references
on public or generated data before expensive bank experiments. Changes in
source, target, covariance law or delivery trigger re-review. The two reference
decisions settle the proposal's present reuse choice while granting no native
integration capability or general sensitivity guarantee.
